\documentclass[UTF8,11pt]{ctexart}
\usepackage[a4paper,margin=22mm]{geometry}
\usepackage{amsmath,amssymb,bm,graphicx,adjustbox,fvextra,longtable,array}
\xeCJKDeclareCharClass{CJK}{"2160 -> "217F, "2460 -> "249B, "2200 -> "22FF}
\newcommand{\fitmath}[1]{\begin{adjustbox}{max width=\linewidth}$\displaystyle #1$\end{adjustbox}}
\setlength{\parindent}{0pt}\setlength{\parskip}{.65em}\renewcommand{\arraystretch}{1.3}
\sloppy\allowdisplaybreaks
\begin{document}
\section*{2010年考研政治答案解析}
\subsection*{2010 年思想政治理论试题答案及解析}

\subsection*{一、单项选择题（每小题1分，共16分）}

\subsection*{1.【标准答案】A}

本题考查马克思主义追求的根本价值目标。本题只要抓住题干的核心意思即可选择出最佳答案。题干恩格斯认为“未来社会主义纪元的基本思想”是“每个人自由而全面的发展”，这表明，马克思主义追求的根本价值目标就是实现人的自由而全面的发展，故 A正确。B是以西方“抽象人性论”为基础的一种社会理想；C是自古就有的空想主义的观点和主张；D是资产阶级在反对封建制度时提出的主张。BCD均不符合题意。


\subsection*{2.【标准答案】C}

本题考查量变质变规律。题干中的“坚持”是一个量的积累过程，即量变，而“胜利”说明事物的根本性质发生了变化即质变，“坚持就是胜利”表明了量变到质变的过程。故C正确。A与题意无关，必然性是指事物联系和发展过程中一定要发生、确定不移的性质和趋势。偶然性是指事物联系和发展过程中并非确定发生的性质和趋势。B与题意无关，肯定因素是维持现存事物存在的因素，否定因素是促使现存事物灭亡的因素，否定是事物的自我否定，而题干中促使“岩石”变化的并非“岩石”本身，而是“溪水”这一外部力量。D干扰性最大，因果联系强调的是“引起”和“被引起”的关系，“坚持就是胜利”这句话主要目的不是强调因为“坚持”所以“胜利”，而在于强调“坚持”这一量的积累过程，故就题干的要求而言，D不及C准确。


\subsection*{3.【标准答案】D}

本题考查人与自然的关系。自然界是人类社会形成的自然基础，人类通过实践影响和制约自然界，不断改变自然界。如果人类用正确的实践方式改造自然，会实现“人类同自然的和解”，如果用错误的实践方式改造自然，最终会遭受大自然的“报复”。气候变暖是人类实践方式不当，破坏生态平衡的结果，由此引发的警示是人与自然必须建立和谐的关系，人类要思考自身活动对环境的影响，尽量避免或减少破坏生态环境，故D正确。A表述错误，一切社会关系的核心是经济关系；B表述错误，生态失衡并非是自然界自身周期演化不可逆转的趋势；C表述错误，生产力才是推动人类社会发展的根本决定力量。


\subsection*{4.【标准答案】C}

本题考查货币转化为资本。资本是能够带来剩余价值的价值，货币本身并不是资本，只有当货币在运动中发生了价值的增殖，带来剩余价值时，货币才能转化为资本。因此，资本家必须购买到一种特殊的商品，这种商品具有创造价值和剩余价值的能力，这种特殊的商品就是劳动力。故C正确。A表述错误，资本家购买的是劳动力这种商品的使用价值和价值。B表述正确，但仅强调了劳动力商品的属性，如果资本家不用货币去购买这种劳动力商品，货币仍然不能转化为资本；D表述错误，劳动力自身的价值在消费过程中不是转移到新的商品中去，而是由工人的劳动再生产出来。


\subsection*{5.【标准答案】A}

本题考查十一届六中全会对我国社会主要矛盾的规范性表述。首先，由题干可知，这一主要矛盾包括“人民日益增长的物质文化需要”和“落后的社会生产”两个方面。前者是矛盾的次要方面，后者是矛盾的主要方面，故可排除D。其次，“落后的社会生产”主要是指生产力的落后，此外，也包括生产力的组织、经营和管理方式的落后。故 A为最佳答案。B虽然也有“生产力”三个字，但其回答的并非主要矛盾的内容，而是解决这一主要矛盾的要求。C从逻辑关系上讲，属于“生产力落后”的具体表现，不及 A全面。


\subsection*{6.【标准答案】D}

本题考查发展的根本目的。本题一定要注意审题，题干问的是“根本目的”。中国解决所有问题的关键要靠自己发展，发展的根本目的是使人民共享发展成果，实现共同富裕。故D正确。ABC 虽然均是发展的目的，但不是“根本”目的，不符合题意。


\subsection*{7.【标准答案】A}

本题考查建立资源节约型社会的核心。资源节约型社会，是指以能源资源高效率利用的方式进行生产、以节


约方式进行消费为根本特征的社会，核心是节约使用能源资源和提高能源资源利用效率，故 A 正确。B是建立资源节约型社会的重要举措，但不是核心。C“限制”一词表述错误，应改为“合理开发和利用”。D是建立资源节约型社会的重要途径，并非核心。


\subsection*{8.【标准答案】C}

本题考查保障人民基本文化权益的主要途径。发展公益性文化事业，建立政府主导的公共文化体系是保障人民基本文化权益的主要途径，C正确。A 是建设社会主义文化的总体目标。B是深化文化体制改革的方针。D是发展经营性文化的重要途径。ABD均与题意不符。


\subsection*{9.【标准答案】B}

本题考查马克思主义在中国的传播。李大钊在中国大地上率先举起马克思主义的大旗，是近代中国第一个马克思主义者。他在《法俄革命之比较观》中第一次向中国人民阐述了十月革命的性质，后来又写了《庶民的胜利》、《布尔什维主义的胜利》、《我的马克思主义观》，对马克思的理论作了比较系统的介绍，B正确。


\subsection*{10.【标准答案】B}

本题考查社会主义道路的初步探索。1956 年《论十大关系》的发表，标志着以毛泽东为主要代表的中国共产党人开始探索中国自己的社会主义建设道路，B正确。A错在“实现了”，《论十大关系》的发表只是开始了马克思主义与中国实际的“第二次结合”。C的表述是毛泽东在1960年才提出来的，与题意不符。D错在“已经”一词，“突破社会主义苏联模式的束缚”是在改革开放后才实现的。


\subsection*{11.【标准答案】B}

本题考查个人品德中的德才关系。题干这句话表达的是“道德（品格）比才智更重要”，讲的是“德”和“才”的关系。B强调“才是德的辅助，德是才的统帅”，说的也是“德”和“才”的关系，故B正确。ACD 均不符合题意。A表达的是“路程虽远，不走就不能到达目的地；事情虽小，不做就不能成功”，强调实践的重要性；C强调“礼”是人的立身之本；D强调能够分辨是非得失，是有智慧的表现。


\subsection*{12.【标准答案】C}

本题考查民族精神的内涵。题干中提及的传说都反映了中国人民不屈不挠的自强不息精神和勇于追求、实现梦想的执着精神，故C正确。上述传说并未体现“勤劳勇敢”、“团结统一”、“爱好和平”，ABD不符合题意。


\subsection*{13.【标准答案】B}

本题考查公民基本道德规范的主要内容。“诚实守信”既是我国公民道德建设的重点，也是中华民族的传统美德，也是职业生活中的基本道德规范，还是社会主义核心价值观的一条重要准则，故选B。ACD 虽然都是公民基本道德规范的内容，但不是重点。


\subsection*{14.【标准答案】A}

本题考查依法治国的根本要求。1978年，邓小平在中央工作会议上的讲话中明确提出：“有法可依、有法必依、执法必严、违法必究”，这十六字方针是依法治国的根本要求，故 A 正确。B是我国政治体制改革的主要任务；C是公平正义观念的基本内容；D讲的是依法治国的基本含义，而非根本要求。


\subsection*{15.【标准答案】D}

\subsection*{16.【标准答案】D}

\subsection*{二、多项选择题（每小题2分，共34分。多选或少选均不得分）}

\subsection*{17.【标准答案】ACD}

本题考查意识的本质和起源。\textcircled{1}从意识的本质来看，意识是物质世界的主观映象。意识的内容来源于人们的实际生活，并随着社会生活的变化而变化发展。《新华字典》所有语词的含义都是对客观世界的能动反映，并且随着人们社会生活的变化而变化。故 AC正确。\textcircled{2}从意识的起源来看，意识不仅是自然界长期发展的产物，而且是社会历史的产物。意识的社会性，不仅表现在劳动为意识的产生和发展提供了客观需要和可能，而且也表现在意识必须借助于语言这一物质外壳才能固定下来、表达出来。《新华字典》以及其他各种词典的发行并一版再版，就是借助语言中的一种形式——纸质文字进行文化传承和思想交流。故D正确。B颠倒了意识与“词语含义”之间的关系。语言是人类意识的物质外壳，是意识的表达形式。人类的社会生活决定了意识的内容并进而规定了词语的表达形式。


\subsection*{18.【标准答案】CD}

本题考查对科学技术的理解。本题是一道灵活的阅读理解题。首先，第1句话总体上强调了促使科技发展的重大原因是经济危机，也即社会实践的需要，故C正确。其次，第2、3句话指出了科技创新带来的重大影响，也即推动世界实现了从电气时代到电子时代的跨越式发展，故D正确。A 表述错误，摆脱社会危机的根本出路是解决社会基本矛盾；B表述错误，社会形态更替的根本标志是社会制度特别是经济制度的根本改变。


\subsection*{19.【标准答案】ACD}

本题考查资本主义政治制度的本质。资本主义政治制度是建立在资本主义生产资料私有制基础之上的，是为资产阶级服务，是服从于资产阶级进行统治和压迫需要的政治工具。题干将狐狸比喻资产阶级，以平底盘子比喻资本主义社会的“平等”（形式上的平等），以仙鹤实际上喝不到平底盘子中的鱼汤比喻资本主义社会“平等”的实质（事实上的不平等），意在说明：\textcircled{1}资本主义社会的平等是建立在财产不平等——资本主义生产资料私有制基础之上的，是形式上的平等，并不是事实上的平等。因此，B错误，D正确；\textcircled{2}资产阶级宣布“法律面前人人平等”，实质是将这种不平等合法化，以法律名义上的平等掩盖着事实上的不平等。因此，AC正确。


\subsection*{20.【标准答案】BCD}

本题考查社会形态更替的前进性与曲折性。社会形态的前进性是指原始社会、奴隶社会、封建社会、资本主义社会、共产主义社会五种社会形态依次演进的基本趋势，曲折性是指社会前进过程中所出现的反复、停滞和倒退现象。“历史终结论”是把社会主义在一些国家的暂时失败看作资本主义的永久胜利，这种理论的破产说明：\textcircled{1}任何事物的发展都不是一帆风顺，人类历史发展是前进性与曲折性的统一，B正确；\textcircled{2}社会主义在一些国家遭到挫败，并不等于整个社会主义事业的失败，C正确；\textcircled{3}人们对社会发展某个阶段的认识不能代替对社会发展的整个过程的认识，D正确。A表述错误，社会规律不像自然规律一样是作为一种盲目的无意识的力量起作用，而是通过抱有一定目的和意图的人的有意识的行动实现的。


\subsection*{21.【标准答案】ABC}

本题考查马克思主义中国化的原因。1938年，毛泽东在党的六届六中全会上作的题为《论新阶段》的政治报告中，第一次提出了“马克思主义中国化”的命题。马克思主义之所以能够中国化，是因为：\textcircled{1}这是马克思主义理论的内在要求，故 A正确；\textcircled{2}这是解决中国实际问题的需要，故C正确；\textcircled{3}马克思主义与中华民族优秀文化具有相融性，故B正确；D表述错误，马克思主义只能为中国的革命、建设和改革提供方法论指导，而不能提供现实的发展模式。


\subsection*{22.【标准答案】ABD}

本题考查过渡时期总路线提出的原因和条件。具体包括：\textcircled{1}到1952年底，一方面，我国的民主革命遗留任务已经完成，国内的阶级关系和主要矛盾发生了深刻的变化。这说明，明确提出向社会主义过渡的任务已经成为必要的了。故D正确。另一方面，随着国民经济的恢复和初步发展，中国社会的经济成分（即生产关系）发生了重要变化。社会主义的国营经济在国家经济生活中实际上已经居于强大的地位，并且明显地表现出对于其他经济成分的优越性。这不仅为党提出的过渡时期的总路线奠定了强大的物质基础，同时，也充分说明对资本主义经济进行限制和改造的必要性。故 A正确；\textcircled{2}历史和现实表明，中国的工业化必须坚持社会主义方向。故B正确。C错在“取得了”三个字，当时我国的国民经济仅是得到了恢复和发展，不可能已经取得了重大成就。我国工业化建设取得重大成就是在第一个五年计划期间。


\subsection*{23.【标准答案】ABCD}

本题考查人民代表大会制度。人民代表大会制度是我国的政体。这一制度是中国共产党把马克思主义与中国实际相结合的伟大创造，是近代以来中国社会发展的必然选择，是中国共产党带领全国人民长期奋斗的重要成果，是全国各族人民的共同利益和共同愿望的反映，故 ABCD正确。


\subsection*{24.【标准答案】BCD}

本题考查处理民族关系的原则。具体包括：维护祖国统一；反对民族分裂；坚持民族平等、民族团结、各民族共同繁荣。故 BCD 正确。A是我国解决民族问题的基本政策，与题意不符。


考妍】后台发送【PDF】新


\subsection*{25.【标准答案】ACD}

本题考查走和平发展道路的依据。中国走和平发展道路：\textcircled{1}是基于中国文化传统的必然选择。中华民族是热爱和平的民族，和合文化是中国追求和平发展的文化基础。中国在近代历史上是个遭受帝国主义欺压的半殖民地半封建国家，中国人民饱经战乱和穷困之苦。中国人民从悲惨遭遇中深刻体会到，和平十分珍贵；\textcircled{2}是基于中国国情的必然选择。社会主义的国家性质决定中国外交的宗旨和首要目标必须是和平与发展；\textcircled{3}是基于当今世界发展潮流的必然选择。和平与发展是当今时代的主题，要和平、促发展、谋合作是各国人民的共同追求，是世界各国人民的共同愿望。一定要注意审题，本题问的是中国走和平发展道路的一个原因一—历史文化传统，只能选 ACD，不能选B。但如果题干问“中国走和平发展道路的依据”则选 ABCD。


\subsection*{26.【标准答案】ABD}

本题考查洋务运动失败的主要原因。具体包括：\textcircled{1}指导思想的封建性。洋务运动的指导思想是“中学为体，西学为用”即在封建主义思想的指导下，在维持封建的上层建筑、经济基础的条件下发展一些近代企业，为维护清王朝的封建统治而服务，这就决定了它必然失败；\textcircled{2}对外国具有依赖性。洋务派兴办的企业一切依赖外国，企图依赖外国达到“自强”、“求富”的目的，这无异于与虎谋皮；\textcircled{3}洋务企业的管理具有腐朽性。洋务派创办的一些新式企业虽然具有一定的资本主义性质，但其管理基本上仍是封建衙门式的。故 ABD 正确。C对洋务运动有一定影响，但不是洋务运动失败的主要原因。


\subsection*{27.【标准答案】ABCD}

本题考查农村包围城市、武装夺取政权道路的依据。具体包括：\textcircled{1}近代中国是一个半殖民地半封建的社会，内无民主制度，外无民族独立，这决定了中国革命斗争形式只能是武装夺取政权，故 A正确；\textcircled{2}近代中国农民占全国人口绝大多数，是无产阶级可靠的同盟军和革命主力军，无产阶级要想夺取革命胜利就必须把工作重点放在农村，获得农民支持，故 B正确；\textcircled{3}中国革命的敌人建立了庞大的军队，并长期占据着中心城市，农村是其统治的薄弱环节，无产阶级政党必须将工作重心放到敌人统治力量薄弱的农村去，故C正确；\textcircled{4}近代中国经济政治发展的不平衡，这为农村建立革命根据地的存在和发展提供了客观的政治经济条件，故D正确。


\subsection*{28.【标准答案】ABD}

本题考查抗日民族统一战线的原则和方针。主要有：\textcircled{1}坚持既联合又斗争的原则；\textcircled{2}采取发展进步势力，争取中间势力，孤立顽固势力的总方针；\textcircled{3}同顽固派斗争坚持有理、有利、有节的原则，故 ABD 正确。C是抗战胜利后毛泽东提出的处理国内阶级斗争的对策。


\subsection*{29.【标准答案】ACD}

本题考查解放战争中第二条战线形成的原因。解放战争中第二条战线，是相对于第一条战线（即军事战线，1946年6月内战全面爆发后形成）而言的，是在国民党统治区形成的以学生运动为先导的人民民主运动。形成的原因包括：\textcircled{1}国民党政府专制独裁、官员贪污腐败；\textcircled{2}国民党违背人民意愿，顽固坚持内战政策；\textcircled{3}国民党政府苛捐杂税，滥发纸币等导致国统区爆发严重经济危机，故 ACD 正确。B是第一条战线的内容，与题意无关。


\subsection*{30.【标准答案】ABCD}

本题考查新时期的爱国主义。这是一道材料考点结合题，解答这类题需同时满足两个原则：一是看选项本身表述是否正确，二是看其是否与题意匹配。材料的主旨是“爱国主义”。这对我们的启示有：\textcircled{1}科学知识是无国界的，但科学事业的发展和科学家的命运都与自己的祖国有密切的关系，故 A正确；\textcircled{2}个人理想如果离开国家和民族就会成为无源之水，无本之木，故B正确；\textcircled{3}在当代中国，爱国主义首先体现在对社会主义中国的热爱上，爱国主义与爱社会主义的统一是中国历史发展的必然结果，故C正确；\textcircled{4}爱国既需要情感的基础，也需要理性的认识，更需要实际的行动，故D正确。


\subsection*{31.【标准答案】BD}

本题考查我国公民基本权利中的政治权利。我国公民的基本权利主要包括：政治权利、人身权利、财产权利、社会经济权利、宗教信仰及文化权利等。政治权利主要包括选举权、被选举权、表达权（政治表达的自由）、民主管理权、监督权等，故 BD正确。人身权利包括生命健康权、人身自由权、人格尊严权、住宅安全


权、通信自由权等，A属于人身权利，故不选；宗教信仰及文化权利包括宗教信仰自由、文化教育权等，C属于宗教信仰及文化权利。


\subsection*{32.【标准答案】ABCD}

\subsection*{33.【标准答案】ABC}

\subsection*{三、分析题 （每小题10分，共50分）}

\subsection*{34.【标准答案】}

（1）人类的认识和实践是一个复杂的过程。由于主客观条件的限制和制约，任何人发生错误都是难免的。（3分）


（2）真理与谬误、成功与失败是对立统一的关系，它们互相包含又能在一定的条件下互相转化。梅兰芳对剧情本身有着深刻的理解，自身具有深厚的艺术实践功力，所以能化险为夷，变失败为成功。（4分）


（3）我们要正视失败和错误，认真总结经验教训，根据实际情况采取恰当的措施和办法并加以应对，促成事物朝有利的方向转化。（3分）


\subsection*{35.【标准答案】}

（1）社会建设与人民幸福息息相关，社会建设是中国特色社会主义事业总体布局的重要组成部分。经济建设是基础，经济建设的最终目的是为了改善民生，提高人民的生活水平。（2分）加快推进以改善民生为重点的社会建设，有利于解决人民群众最关心、最直接、最现实的利益问题，是贯彻落实科学发展观的重要内容，关系到巩固党执政的社会基础。（3分）


（2）把以改善民生为重点的社会建设贯彻到中国特色社会主义建设的全过程，并在社会主义建设的各个方面体现出来。（2分）优先发展教育，实现教育公平；实现扩大就业的发展战略，促进以创业带动就业；深化收入分配制度改革，增加城乡居民收入；加快建立覆盖城乡居民的社会保障体系，保障人民基本生活；建立基本医疗卫生制度，提高全民健康水平；完善社会管理，维护社会安定团结。（3分）


\subsection*{36.【标准答案】}

（1）孙中山立志救亡图存，振兴中华，领导辛亥革命，结束中国的君主专制制度，建立资产阶级共和国，但是民主共和国的理想并没有实现。（2分）中国共产党人继承孙中山的遗志，开辟了中国革命的正确道路和新的发展方向，领导全国人民进行艰苦卓绝的斗争，终于完成反帝反封建的民主革命任务，建立了中华人民共和国。（3分）


（2）国家统一基本完成，对外获得民族独立；人民民主专政国家政权建立，人民当家做主；中国共产党成为执政党。这些为社会主义基本制度的建立和当代中国一切发展进步奠定了根本政治前提。（3分）新中国成立以来，中国共产党团结带领全国各族人民在革命、建设、改革的伟大实践中，取得了举世瞩目的伟大成就，贫穷落后的中国变成了一个初步繁荣昌盛，充满生机和活力的社会主义国家。（2分）


\subsection*{37.【标准答案】}

（1）文明出行是现代社会公共生活的重要内容，随着社会的进步，公共生活领域的范围逐渐扩大。（2分）公共生活需要公共秩序。维护公共秩序对经济社会健康发展、保障人民生活质量与安全具有重要意义。（2分）建立和维护社会秩序需要道德和法律两种手段。两者发挥作用的方式有所不同，但互为补充，相辅相成。（2分）


（2）构建文明的公共生活秩序，需要增强社会公德意识和法律意识，养成遵守社会公德和遵纪守法的良好行为习惯，学习和把握公共生活中的道德与法律规范，提升自身文明素质。（4分）


\subsection*{38.【标准答案】}

（1）在经济全球化背景下，世界的生产贸易活动紧密联系在一起，相互依存，相互开放，中国制造即世界制造，中国离不开世界，世界需要中国；改革开放以来，中国与世界相互渗透，你中有我，我中有你，“中国制造”对世界的发展不是威胁，而是贡献。中国获益，世界也获益。（5分）


（2）我国还处在世界产业链的较低位置，走自主创新之路，让“中国制造”尽快成为“中国创造”，是时代发展提出的迫切要求；当代世界，综合国力的竞争深刻表现为一场世界范围内“创新力”的竞争。创新，才能提高我国产品的国际竞争力，才能促进我国调整产业结构，转变发展方式，建立创新型国家，进而提升我国的软实力。（5分）

\end{document}
